\sgaStableTarget{sga6.e10.sec.7.15}\subsection*{7.15. Varying @@M0@@.}
Let @@M1@@ be a second trait, and let
@@M2@@
be a morphism of traits, i.e., a dominant local morphism (i.e., carrying the closed point to the closed point and the generic point to the generic point). Denote by @@M3@@ (resp. @@M4@@) the closed point (resp. the generic point) of @@M5@@, and set @@M6@@. Suppose that @@M7@@ is noetherian (e.g., @@M8@@ of finite type over @@M9@@). There is a commutative diagram
@@M10@@
whose oblique arrows are flat morphisms. One can therefore take the inverse-image homomorphisms for @@M11@@ along all the arrows of the diagram (IV 2.12). Applying definition 7.3, one finds that the following diagram
\sgaStableTarget{sga6.e10.eq.7-15-1}@@M12@@
is commutative, where the vertical arrows are the inverse-image homomorphisms. One immediately obtains an analogous commutative diagram with @@M13@@ in place of @@M14@@; and when @@M15@@ and @@M16@@ are regular, one can also write @@M17@@ and @@M18@@ in place of @@M19@@ and @@M20@@.

